\scanpage{93}\mspage{37}
% L1-2
\deleted{\unsure{D'aut}} Tous les \inserted{quatre} foncteurs $i,j,k,\Phi_0$ sont compatibles aux localiseurs, donc passent aux localisés
% L3-5
\diagram{\begin{tikzcd}[ampersand replacement=\&,column sep=small]
 \mathrm{HOT}_W(S)\arrow[dr,"{\bar j}"']\&\& \mathrm{HOT}^{?}_W(S)\arrow[ll,"{\bar i}"']\arrow[dl,"{\bar k}"']\\
 \&\mathrm{HOT}^{\mathrm{lc}}_W(S)\arrow[ur,bend right=32,"{\bar{\unsure{\varphi|\psi}}}"']\&
\end{tikzcd}}
,
% L6-9
où l'exposant $?$ dans la notation signifie que je ne sais interpréter cette catégorie, sauf si je prouve que $\bar k$ est une équivalence\unsure{;|,} et l'exposant $\mathrm{lc}$ dans la catégorie du bas
% L10-14
doit suggérer \deleted{\unsure{ma présomption}} \inserted{l'intuition} qu'il s'agit \deleted{\unsure{d'} la \unsure{sous-c}} en tous cas d'une sorte de ``types d'homotopie relatifs'' (dans un sens \unsure{modifié}) \uline{localement constants} sur $S$. Le problème est de comprendre ce diagramme.

% L15-21
La conjecture naturelle est

1\textsuperscript{o}) $\bar k$ est une équivalence. (Ce qui impliquera que $\bar\varphi$ est l'équivalence quasi-inverse\unsure{)}, que \deleted{2\textsuperscript{o})} $\bar k$, \inserted{$\bar\varphi$} permettent \unsure{essentiellement} d'identifier $\mathrm{HOT}^{?}_W$ et $\mathrm{HOT}^{\mathrm{lc}}_W$, de sorte que $\bar i$ et $\bar j$ deviennent des foncteurs entre $\mathrm{HOT}_W$ et $\mathrm{HOT}^{\unsure{\mathrm{loc}|\mathrm{lc}}}_W$.)

% L22-24
2\textsuperscript{o}) $\bar i$ est pleinement fidèle (NB l'image essentielle est bien comprise de \unsure{toute façons|toute façon}.)
% archive sheet number: 92
